\documentclass[12pt,a4paper]{article}

\usepackage{cmap}
\usepackage[T2A]{fontenc}
\usepackage[utf8]{inputenc}
\usepackage[english,russian]{babel}
\usepackage[a4paper,left=2.5cm,right=2cm,top=2cm,bottom=2cm]{geometry}
\usepackage{booktabs}
\usepackage{array}
\usepackage{enumitem}
\usepackage{amsmath}
\usepackage[table]{xcolor}
\usepackage{hyperref}

\definecolor{candy}{RGB}{217,79,112}
\definecolor{choco}{RGB}{74,44,23}

\hypersetup{
  colorlinks=true,
  linkcolor=candy,
  urlcolor=candy,
  pdftitle={Конфетная фабрика — руководство по балансировке},
  pdfauthor={Конфетная фабрика}
}

\setlength{\parindent}{0pt}
\setlength{\parskip}{6pt}

\newcommand{\A}{A_0}
\newcommand{\vv}{v_0}

\title{\textbf{\textcolor{choco}{Конфетная фабрика}}\\[4pt]
\large руководство ведущего по балансировке параметров}
\date{}

\begin{document}
\maketitle
\vspace{-2.5em}

Это руководство объясняет, \emph{как устроена экономика игры} и какие
значения ставить в форме создания игры, чтобы соревнование было честным,
азартным и не разваливалось ни для сильных, ни для слабых команд.
Все выводы аналитические, в конце --- готовые пресеты и чек-лист.

\section{Модель экономики}

Обозначения (для уровня сложности $r$):

\begin{center}
\begin{tabular}{llll}
\toprule
Параметр в форме & Символ & По умолчанию ($n=3$) \\
\midrule
Длительность, сек        & $T$      & 5100 (85 мин) \\
Стартовые запасы         & $\A$     & 20\,000 \\
Стартовая скорость       & $\vv$    & 15 \\
Цена задачи              & $C_r$    & 12\,000 / 12\,000 / 12\,000 \\
Цена подсказки (теста)   & $H_r$    & 3\,000 / 7\,000 / 10\,000 \\
Нагрузка от задачи       & $L_r$    & 2 / 2 / 2 \\
Бонус запасы             & $B_r$    & 12\,000 / 25\,000 / 50\,000 \\
Бонус скорость           & $S_r$    & 4 / 7 / 11 \\
\bottomrule
\end{tabular}
\end{center}

Механика: каждую секунду склад пополняется на текущую скорость.
Покупка задачи: $-C_r$ к запасам, $-L_r$ к скорости. Решение:
$+B_r$ к запасам, скорость возвращает $L_r$ и получает $+S_r$ сверху.
Побеждает наибольший склад в момент $T$.

\subsection{Главная формула}

Пусть задача уровня $r$ куплена в момент $t$ и решена за время $\tau$
(секунд). Тогда её суммарный вклад в итоговый счёт:
\begin{equation}
\Delta A \;=\; \underbrace{B_r - C_r}_{\text{мгновенная прибыль}}
\;-\; \underbrace{L_r\,\tau}_{\text{потери на время решения}}
\;+\; \underbrace{S_r\,\bigl(T - t - \tau\bigr)}_{\text{разгон до конца игры}}.
\label{eq:profit}
\end{equation}

Из \eqref{eq:profit} следуют все балансные свойства игры:

\begin{itemize}[itemsep=2pt]
\item \textbf{Скорость --- главный приз.} Слагаемое $S_r(T{-}t{-}\tau)$
растёт с каждой минутой, оставшейся до конца: решать \emph{рано} гораздо
выгоднее, чем решать \emph{много} в конце.
\item \textbf{Цена минуты промедления.} Решение на минуту раньше дороже
на $(L_r + S_r)\cdot 60$ конфет: при значениях по умолчанию это
\textbf{360 / 540 / 780} конфет за минуту по уровням. Хотите подстегнуть
темп --- увеличивайте $S_r$; хотите спокойную игру --- уменьшайте.
\item \textbf{Порог окупаемости.} Задача, решённая, когда до конца
остаётся $R$ секунд, окупается при
$R > R^{*} = \dfrac{C_r + L_r\tau - B_r}{S_r}$.
Если $B_r \ge C_r + L_r\tau$, задача выгодна \emph{всегда} --- даже
решённая на последней секунде.
\end{itemize}

Порог $R^{*}$ при параметрах по умолчанию (в минутах остатка):

\begin{center}
\begin{tabular}{lccc}
\toprule
Время решения $\tau$ & 15 мин & 30 мин & 45 мин \\
\midrule
Уровень 1 & 7{,}5 мин & 15 мин & 22{,}5 мин \\
Уровень 2 & всегда выгодна & всегда & всегда \\
Уровень 3 & всегда выгодна & всегда & всегда \\
\bottomrule
\end{tabular}
\end{center}

То есть по умолчанию средние и сложные задачи стоит решать до самого
финала, а лёгкие перестают окупаться только в последние минуты ---
это осознанный выбор: никто не сидит без дела в конце.

\subsection{Сколько это в конфетах: опорные сценарии}

Пассивная команда (ничего не покупает) финиширует с
$\A + \vv T = 96\,500$. Вклад одной задачи, купленной на 10-й минуте
(формула \eqref{eq:profit}):

\begin{center}
\begin{tabular}{lccccc}
\toprule
Решена за & 10 мин & 20 мин & 30 мин & 45 мин & 60 мин \\
\midrule
Уровень 1 & $+14\,400$ & $+10\,800$ & $+7\,200$ & $+1\,800$ & $-3\,600$ \\
Уровень 2 & $+39\,100$ & $+33\,700$ & $+28\,300$ & $+20\,200$ & $+12\,100$ \\
Уровень 3 & $+79\,700$ & $+71\,900$ & $+64\,100$ & $+52\,400$ & $+40\,700$ \\
\bottomrule
\end{tabular}
\end{center}

Типичные траектории целых команд (последовательные покупки, без тестов):

\begin{center}
\begin{tabular}{lrr}
\toprule
Команда & Финиш & К пассиву \\
\midrule
Пассивная                                  & 96\,500  & $1{,}0\times$ \\
Слабая: 3 лёгкие по 18 мин                 & 120\,260 & $1{,}25\times$ \\
Средняя: 2 лёгкие + 2 средние              & 165\,760 & $1{,}7\times$ \\
Сильная: по 2 задачи каждого уровня        & 283\,280 & $2{,}9\times$ \\
\bottomrule
\end{tabular}
\end{center}

Это «здоровый» профиль: активность всегда вознаграждается, разрыв
между полюсами --- примерно втрое, интрига сохраняется. При настройке
своих параметров стремитесь к похожей картине (как посчитать --- ниже).

\section{Рецепт настройки за 6 шагов}

Единственное, что нужно оценить самостоятельно, --- \emph{сколько минут
медианная команда решает задачу каждого уровня} ($\tau_1 < \tau_2 < \dots$).
Остальное выводится по правилам ниже.

\begin{enumerate}[leftmargin=1.6em,itemsep=3pt]

\item \textbf{Длительность $T$ и размер доски.}
90 минут --- классика при $n=3$ (9 задач). 120 минут --- либо те же
9 задач, но сложнее, либо $n=4$ (16 задач). Ориентир загрузки: сумма
$n\cdot(\tau_1{+}\dots{+}\tau_n)$ человеко-минут должна быть
\emph{близка или чуть больше} командного бюджета времени
($\approx 3\cdot T$ человеко-минут для команды из трёх), тогда лучшая
команда соберёт почти всю доску лишь к самому концу. По умолчанию:
$3\cdot(13{+}25{+}45)\approx 250$ против бюджета $255$ --- впритык,
как и задумано.

\item \textbf{Скорость и нагрузка.} $\vv$ --- любое круглое число,
15 подходит всем. Нагрузку держите в диапазоне
$L_r \approx \vv/8 \dots \vv/6$ (для $\vv=15$ это~2): тогда команда
может вести 3--4 задачи параллельно, не обнуляя конвейер,
а покупка всё же ощущается.

\item \textbf{Цена задачи.} Единая для всех уровней (различия ---
в бонусах, так проще принимать решения):
\[
C \approx (0{,}12\dots0{,}18)\cdot \vv\,T .
\]
Для $T=85$ мин: $\vv T = 76\,500$, отсюда $C \approx 12\,000$.
Смысл: задача стоит 10--15 минут работы конвейера --- решение о покупке
значимое, но не страшное.

\item \textbf{Стартовые запасы.} $\A \approx (1{,}5\dots2)\cdot C$:
одну задачу можно купить сразу, вторую --- через несколько минут.
Меньше $C$ не ставьте никогда (мёртвый старт).

\item \textbf{Бонусы --- сердце баланса.}
\[
B_1 \approx C, \qquad B_n \approx (3\dots4)\cdot C,
\qquad \text{промежуточные — геометрически.}
\]
$B_1{=}C$ делает лёгкий уровень «безрисковым обучением», а премия
за сложность растёт с риском. Бонус скорости --- пропорционально
времени решения ($S_r \sim \tau_r$), с масштабом
\[
S_r \cdot \tfrac{T}{2} \approx (0{,}7 \dots 2{,}3)\cdot C
\quad\text{(от лёгкого уровня к сложному)}.
\]
Для умолчаний: $4\cdot2550=10\,200\approx0{,}85C$ и
$11\cdot2550=28\,050\approx2{,}3C$. Проверка сверху: потолок скорости
при полностью решённой доске $\vv + n\sum_r S_r$ должен быть
в $4\dots6$ раз больше $\vv$ (по умолчанию $81 = 5{,}4\,\vv$).
Больше --- «снежный ком», меньше --- решать задачи слабо мотивирует.

\item \textbf{Цена подсказки.} От четверти до~80\,\% цены задачи,
растёт с уровнем: $H_r \approx (0{,}25 / 0{,}6 / 0{,}8)\cdot C$.
Дешевле --- команды скупают тесты вместо размышлений, дороже ---
кнопкой никто не пользуется.

\end{enumerate}

\subsection{Как связать уровни с реальными задачами}

Ориентиры для informatics при 90-минутной игре и школьной аудитории:

\begin{center}
\begin{tabular}{llll}
\toprule
Уровень & Ожидаемое $\tau$ & Кто решает & Подбор задач \\
\midrule
1 (Easy)   & 10--15 мин & все команды      & «разогрев», 1--2 идеи, простая реализация \\
2 (Middle) & 20--35 мин & большинство      & одна содержательная идея + аккуратный код \\
3 (Hard)   & 40--60 мин & 1--2 лучших      & полноценная олимпиадная задача \\
\bottomrule
\end{tabular}
\end{center}

Если реальные $\tau$ заметно отличаются от ориентиров --- пересчитайте
$S_r$ (шаг 5) под свои значения, остальное можно не трогать.

\section{Готовые пресеты}

\newcommand{\preset}[1]{\textbf{#1}}

\begin{center}
\small
\begin{tabular}{lccc}
\toprule
 & \preset{Классика} & \preset{Марафон} & \preset{Разминка} \\
 & 90 мин, $n=3$ & 120 мин, $n=3$ & 60 мин, $n=3$ \\
\midrule
Длительность, мин   & 85       & 115      & 55 \\
Стартовые запасы    & 20\,000  & 24\,000  & 14\,000 \\
Стартовая скорость  & 15       & 15       & 15 \\
Цена задачи         & 12\,000  & 16\,000  & 8\,000 \\
Цена подсказки      & 3/7/10 тыс. & 4/9/13 тыс. & 2/4/6 тыс. \\
Нагрузка            & 2/2/2    & 2/2/2    & 2/2/2 \\
Бонус запасы        & 12/25/50 тыс. & 16/34/64 тыс. & 8/16/28 тыс. \\
Бонус скорость      & 4/7/11   & 4/7/11   & 5/8/12 \\
Задачи ($\tau$, мин) & 12/25/45 & 15/35/60 & 8/15/25 \\
\bottomrule
\end{tabular}
\end{center}

\begin{itemize}[itemsep=2pt]
\item \preset{Классика} --- значения по умолчанию, проверенный баланс.
\item \preset{Марафон} --- та же динамика на 2 часа: деньги
масштабированы вместе с $\vv T$, задачи заметно сложнее. Вариант:
вместо усложнения задач взять $n=4$ (16 задач уровней
$\tau \approx 10/20/35/55$; бонусы запасов 12/22/40/70 тыс.,
скорости 4/6/9/13, подсказки 3/6/9/12 тыс.).
\item \preset{Разминка} --- короткая игра для новичков: всё дешевле,
скоростные бонусы чуть выше (за 55 минут они успевают меньше поработать).
\end{itemize}

\section{Чек-лист перед стартом}

Проверьте свои параметры по пяти инвариантам:

\begin{enumerate}[leftmargin=1.6em,itemsep=2pt]
\item $\A \ge C_{\max}$ --- любая задача доступна с первой секунды
(лучше $\A \approx 1{,}5\,C$).
\item $L_r \le \vv/6$ --- три параллельные задачи не парализуют конвейер.
\item $B_1 \ge C_1$ --- лёгкий уровень безрисковый; $B_n/C_n = 3\dots4$ ---
премия за смелость.
\item Потолок скорости $\vv + n\sum S_r = (4\dots6)\cdot \vv$.
\item Прикиньте по формуле \eqref{eq:profit} финиш «сильной» и «слабой»
траекторий: здоровый разрыв --- $2\dots3\times$, и обе выше пассивной.
\end{enumerate}

\section{Типичные ошибки}

\begin{itemize}[itemsep=2pt]
\item \textbf{«Снежный ком».} Слишком большие $B_n$ и $S_n$: лидер
после первой сложной задачи недосягаем. Симптом: разрыв $>3\times$
к середине игры. Лекарство: снизить $B_n$ до $3C$ и $S_n$ на треть.
\item \textbf{«Зачем решать лёгкие».} $B_1 < C_1$ при малом $S_1$:
во второй половине игры лёгкие задачи убыточны, новички выключаются
из игры. Держите $B_1 \approx C_1$.
\item \textbf{«Паралич».} $L \ge \vv/3$: две купленные задачи почти
останавливают конвейер, команды боятся покупать.
\item \textbf{«Мёртвый старт».} $\A < C$: первые минуты все просто
смотрят на таймер.
\item \textbf{Продление «из жалости».} Помнить: посылка, отправленная
до конца, зачтётся автоматически, даже если тестируется после финиша
(система дособирает решения ещё 15 минут). Продлевать игру ради
«ещё чуть-чуть» не нужно --- для этого есть кнопка продления в админке,
но она сдвигает и порог окупаемости всех купленных задач.
\item \textbf{Разные $\vv$ или $\A$ «для уравнивания команд».}
Технически стартовые условия одинаковы для всех; выравнивайте силу
составов, а не экономику.
\end{itemize}

\vspace{1em}
\begin{center}
\emph{Формула \eqref{eq:profit} + пять инвариантов чек-листа ---
этого достаточно, чтобы собрать свой баланс под любую аудиторию.}
\end{center}

\end{document}
